% Generated draft: requires mathematical and source review.
\documentclass[11pt]{article}
\usepackage[margin=25mm]{geometry}
\usepackage{fontspec}
\setmainfont{Latin Modern Roman}
\usepackage{amsmath,amssymb,mathtools,unicode-math}
\setmathfont{Latin Modern Math}
\usepackage{enumitem,needspace,xurl,hyperref}
\setlength{\parindent}{0pt}
\setlength{\parskip}{5pt}
\setlength{\emergencystretch}{3em}
\newcommand{\N}{\mathbb{N}}
\newcommand{\Z}{\mathbb{Z}}
\newcommand{\Q}{\mathbb{Q}}
\newcommand{\R}{\mathbb{R}}
\newcommand{\C}{\mathbb{C}}
\newcommand{\F}{\mathbb{F}}
\newcommand{\A}{\mathbb{A}}
\newcommand{\PP}{\mathbb{P}}
\newcommand{\E}{\mathbb{E}}
\newcommand{\eps}{\varepsilon}
\newcommand{\dd}{\,\mathrm{d}}
\begin{document}
\section*{UnsolvedMath ID: 20000337}
\textbf{Source problem code:} catalogue entry
\subsection{Definitions and mathematical statement}
Hassett's AIM question asks whether the moduli space of stable surfaces has a virtual fundamental class. In the KSBA interpretation, stable surfaces are projective semi-log-canonical surfaces with ample \(\Q\)-Cartier canonical class. A precise question must fix invariants and distinguish a moduli stack \(\mathcal M\) from its coarse space, as well as the deformation theory, expected dimension, and coefficient theory for a virtual cycle. For example a Behrend--Fantechi perfect obstruction theory is a morphism \(E^\bullet\to L_{\mathcal M}\), with \(E^\bullet\) perfect of amplitude \([-1,0]\), inducing an isomorphism on \(h^0\) and a surjection on \(h^{-1}\); it yields a virtual class of its virtual dimension. The original does not fix these data.

The supplied report says Jiang's arXiv:2206.00575v7 claims a global rational Chow class on a fixed-invariant KSBA stack, by pushing forward the virtual fundamental class of a proper auxiliary stack of lci covers. It explicitly distinguishes this from a perfect obstruction theory on the full KSBA stack, whose natural index-one deformation complex can have higher obstructions. This records the cited version's claim, not an independently verified existence theorem.

The report separately asserts that a universal proper flat lci and Gorenstein family of stable stacky surfaces with finite reduced automorphism groups has a two-term Behrend--Fantechi perfect obstruction theory. Precise base, stack, and universal-family hypotheses are not supplied; this is a restricted lci assertion, not a perfectness statement across the entire non-lci boundary.

Its candidate comparison criterion is the following. Let \(U\subset\mathcal M\) be the intrinsic lci open locus and \(i:B=\mathcal M\setminus U\hookrightarrow\mathcal M\) its closed complement. Suppose two proper auxiliary constructions give classes
\[
\alpha_j=(f_j)_*[\mathcal N_j]^{\mathrm{vir}}\in A_v(\mathcal M)_{\Q}
\quad(j=1,2)
\]
of the same virtual dimension \(v\), and their restrictions agree on \(U\). Then the report asserts
\[
\alpha_1-\alpha_2\in
\operatorname{im}\bigl(i_*:A_v(B)_{\Q}\to A_v(\mathcal M)_{\Q}\bigr).
\]
Consequently \(\alpha_1=\alpha_2\) if \(A_v(B)_{\Q}=0\), in particular if \(\dim B<v\). This is presentation-independence among the compatible constructions agreeing on the lci locus. It does not establish that every proposed auxiliary construction has that agreement, nor does it produce a full-stack obstruction theory. The report requires the rational Chow localization framework in which these restrictions and proper pushforwards are defined.
\subsection{Short English statement}
Ask for virtual classes on stable-surface moduli. The report distinguishes an auxiliary pushforward class from a full perfect obstruction theory and proposes that two constructions agreeing on the lci locus can differ only by a non-lci boundary class.
\subsection{Sources}
\begin{itemize}
\item AIM workshop ``Compact moduli spaces and birational geometry,'' Hassett's question extracted as ``Problem 4. 3.'': \url{https://aimath.org/WWN/birational/birational.pdf}. No PDF was accessed here.
\item Local canonical registry \texttt{aim-algebraic-geometry-notes.json}, \texttt{notes[336]}, tag \texttt{problem}, source item recorded as 4; permanent UnsolvedMath ID \(20000337\), source code \texttt{catalogue entry}. Exact substantive original: ``Does the moduli space of stable surfaces have a virtual fundamental class?'' Original remarks and literature field are empty.
\item Jiang, arXiv:2206.00575v7, \url{https://arxiv.org/abs/2206.00575v7}, as cited by the supplied research report for its auxiliary lci-cover pushforward construction. The report refers to the Behrend--Fantechi obstruction-theory framework without a full citation. No article title or further bibliography is supplied.
\item Supplied research attempt 1 proposes the lci criterion and boundary-support comparison. It is machine generated, not peer reviewed, and requires independent verification before citation.
\end{itemize}
\end{document}
